\documentclass[11pt]{article}
\usepackage[utf8]{inputenc}
\usepackage[T1]{fontenc}
\usepackage{lmodern}
\usepackage[spanish,es-noshorthands,es-nodecimaldot,es-tabla]{babel}
\usepackage[letterpaper,margin=2.2cm,headheight=14pt]{geometry}
\usepackage{amsmath,amssymb,bm}
\usepackage{xcolor}
\usepackage{graphicx}
\usepackage{enumitem}
\usepackage{booktabs,tabularx,array,longtable}
\usepackage[siunitx]{circuitikz}
\usepackage{tikz}
\usetikzlibrary{arrows.meta,calc,decorations.markings,patterns}
\usepackage[most]{tcolorbox}
\usepackage{fancyhdr}
\usepackage{multicol}
\usepackage{microtype}
\usepackage[hidelinks]{hyperref}

\definecolor{azul}{RGB}{20,60,120}
\definecolor{rojo}{RGB}{170,30,30}
\definecolor{verde}{RGB}{20,110,60}
\definecolor{gris}{RGB}{90,90,90}
\definecolor{fondo}{RGB}{244,247,252}

\ctikzset{bipoles/length=1.0cm}

\newcommand{\fem}{\mathcal{E}}
\newcommand{\vE}{\vec{E}}
\newcommand{\vj}{\vec{j}}
\newcommand{\vD}{\vec{D}}
\newcommand{\dS}{\,d\vec{S}}
\newcommand{\dl}{\,d\vec{l}}
\newcommand{\ohm}{\,\Omega}
\newcommand{\uu}[1]{\,\mathrm{#1}}

\tcbset{
  base/.style={enhanced,breakable,boxrule=0.6pt,arc=2pt,left=6pt,right=6pt,top=4pt,bottom=4pt,fonttitle=\bfseries},
}
\newtcolorbox{idea}[1][Idea clave]{base,colback=fondo,colframe=azul,title=#1}
\newtcolorbox{alerta}[1][Error frecuente]{base,colback=red!4,colframe=rojo,title=#1}
\newtcolorbox{metodo}[1][Método]{base,colback=green!4,colframe=verde,title=#1}
\newtcolorbox{profe}[1][Lo que el profesor espera]{base,colback=yellow!6,colframe=orange!70!black,title=#1}

\setlist{itemsep=2pt,topsep=3pt}
\setlength{\parindent}{0pt}
\setlength{\parskip}{4pt}

\pagestyle{fancy}
\fancyhf{}
\renewcommand{\headrulewidth}{0.4pt}
\fancyfoot[C]{\small\thepage}

% Circuitos compartidos entre enunciados y solucionario
% Convención: battery1 con invert dibuja la placa larga (+) en el nodo final del trazo.

\newcommand{\circSunoQdoce}{%
\begin{circuitikz}[scale=0.9,transform shape]
\draw (0,0) to[battery1,invert,l=$\fem_1$] (0,3) to[R,l=$R_1$] (3,3)
      to[R,l_=$R_2$,*-*] (3,0) -- (0,0);
\draw (3,3) to[R,l=$R_3$] (6,3); \draw (3,0) -- (6,0) to[battery1,invert,l_=$\fem_2$] (6,3);
\node[above] at (3,3.25) {$A$}; \node[below] at (3,0) {$B$};
\end{circuitikz}}

\newcommand{\circPuente}[5]{% a-c, a-d, c-b, d-b, c-d
\begin{circuitikz}[scale=0.85,transform shape]
\coordinate (a) at (0,0); \coordinate (b) at (5,0);
\coordinate (c) at (2.5,2); \coordinate (d) at (2.5,-2);
\draw (-1,0) node[left]{$a$} to[short,o-*] (a);
\draw (b) to[short,*-o] (6,0) node[right]{$b$};
\draw (a) to[R,l=#1] (c) to[R,l=#3] (b);
\draw (a) to[R,l_=#2] (d) to[R,l_=#4] (b);
\draw (c) to[R,l=#5,*-*] (d);
\node[above] at (c) {$c$}; \node[below] at (d) {$d$};
\end{circuitikz}}

\newcommand{\circSunoQdiecinueve}{%
\begin{circuitikz}[scale=0.9,transform shape]
\draw (0,0) to[battery1,invert,l={$\fem,\ r$}] (0,3) to[R,l=$R_1$] (3.5,3) to[R,l=$R_2$,*-*] (3.5,0) -- (0,0);
\draw (3.5,3) -- (6,3) to[C,l=$C$] (6,0) -- (3.5,0);
\end{circuitikz}}

\newcommand{\circCuadrado}{%
\begin{tikzpicture}[scale=1.1]
\coordinate (A) at (0,0); \coordinate (B) at (2.5,0); \coordinate (C) at (2.5,2.5); \coordinate (D) at (0,2.5);
\draw[thick] (A)--(B)--(C)--(D)--cycle; \draw[thick] (A)--(C);
\foreach \p/\n/\pos in {A/A/below left,B/B/below right,C/C/above right,D/D/above left}{\fill (\p) circle(2pt); \node[\pos] at (\p){$\n$};}
\node[below] at (1.25,0){$R$}; \node[right] at (2.5,1.25){$R$}; \node[above] at (1.25,2.5){$R$}; \node[left] at (0,1.25){$R$}; \node[above left] at (1.3,1.2){$R$};
\end{tikzpicture}}

\newcommand{\circSdosQdoce}{%
\begin{circuitikz}[scale=0.9,transform shape]
\draw (0,0) to[battery1,invert,l=$\fem_1$] (0,1.6) to[R,l=$R_1$] (0,3.4) -- (3,3.4);
\draw (3,0) to[battery1,invert,l=$\fem_2$] (3,1.6) to[R,l=$R_2$] (3,3.4);
\draw (6,1.6) to[battery1,invert,l_=$\fem_3$] (6,0);
\draw (6,1.6) to[R,l_=$R_3$] (6,3.4) -- (3,3.4);
\draw (0,0) -- (6,0);
\fill (3,3.4) circle(2pt) node[above]{$A$}; \fill (3,0) circle(2pt) node[below]{$B$};
\end{circuitikz}}

\newcommand{\circEscalera}[2]{% serie, paralelo
\begin{circuitikz}[scale=0.85,transform shape]
\draw (0,2) node[left]{$a$} to[short,o-] (0.3,2) to[R,l=#1,-*] (2.3,2) to[R,l=#1,-*] (4.6,2) to[R,l=#1,-*] (6.9,2) -- (7.6,2);
\draw (0,0) node[left]{$b$} to[short,o-] (7.6,0);
\draw (2.3,2) to[R,l=#2,-*] (2.3,0); \draw (4.6,2) to[R,l=#2,-*] (4.6,0); \draw (6.9,2) to[R,l=#2,-*] (6.9,0);
\node at (8.2,1) {$\cdots$};
\end{circuitikz}}

\newcommand{\circInterruptor}{%
\begin{circuitikz}[scale=0.85,transform shape]
\draw (-1,0) node[left]{$a$} to[short,o-*] (0,0) -- (0,1.5) to[R,l=$R$] (2.5,1.5) to[R,l=$2R$] (5,1.5) -- (5,0);
\draw (0,0) -- (0,-1.5) to[R,l_=$2R$] (2.5,-1.5) to[R,l_=$R$] (5,-1.5) -- (5,0);
\draw (5,0) to[short,*-o] (6,0) node[right]{$b$};
\draw (2.5,1.5) to[nos,l=$K$,*-*] (2.5,-1.5);
\node[above] at (2.5,1.6) {$c$}; \node[below] at (2.5,-1.6) {$d$};
\end{circuitikz}}

\newcommand{\circSdosQdiecinueve}{%
\begin{circuitikz}[scale=0.85,transform shape]
\draw (0,0) to[battery1,invert,l=$\fem$] (0,4) -- (6,4);
\draw (3,4) to[R,l=$R_1$,*-*] (3,2) to[R,l=$R_2$,-*] (3,0);
\draw (6,4) to[C,l=$C_1$] (6,2) to[C,l=$C_2$,*-] (6,0) -- (0,0);
\draw (3,2) to[nos,l=$K$] (6,2);
\node[left] at (3,2) {$c$}; \node[right] at (6.1,2) {$d$};
\end{circuitikz}}

\newcommand{\circTetraedro}{%
\begin{tikzpicture}[scale=1.2]
\coordinate (1) at (0,0); \coordinate (2) at (3,0); \coordinate (3) at (1.5,2.6); \coordinate (4) at (1.5,0.9);
\draw[thick] (1)--(2)--(3)--cycle; \draw[thick] (1)--(4)--(2); \draw[thick] (3)--(4);
\foreach \p/\pos in {1/below left,2/below right,3/above,4/below}{\fill (\p) circle(2pt); \node[\pos] at (\p){$\p$};}
\end{tikzpicture}}

\newcommand{\circSdosK}{%
\begin{circuitikz}[scale=0.85,transform shape]
\draw (0,0) node[below]{$G$} to[battery1,invert,l=$\fem_1$,*-] (0,2.5) to[R,l=$R_1$] (0,5) node[above]{$A$};
\draw (0,5) to[R,l=$R_2$,*-*] (8,5) node[above]{$B$};
\draw (0,5) to[R,l_=$R_3$] (4,3.2) node[below left]{$C$};
\draw (4,3.2) to[R,l_=$R_4$,*-] (8,5);
\draw (8,5) to[R,l=$R_5$] (8,0) -- (0,0);
\draw (4,0) to[battery1,invert,l=$\fem_2$,*-] (4,1.6) to[R,l=$R_6$] (4,3.2);
\end{circuitikz}}

\newcommand{\circRueda}{%
\begin{tikzpicture}[scale=1.1]
\foreach \i in {0,...,5}{\coordinate (v\i) at ({60*\i}:1.8);}
\coordinate (O) at (0,0);
\draw[thick] (v0)--(v1)--(v2)--(v3)--(v4)--(v5)--cycle;
\foreach \i in {0,...,5}{\draw[thick] (O)--(v\i); \fill (v\i) circle(2pt);}
\fill (O) circle(2pt); \node[below right] at (O){$O$};
\node[right] at (v0){$1$}; \node[above right] at (v1){$2$}; \node[above left] at (v2){$3$}; \node[left] at (v3){$4$}; \node[below left] at (v4){$5$}; \node[below right] at (v5){$6$};
\end{tikzpicture}}

\newcommand{\circInterruptoresDos}{%
\begin{circuitikz}[scale=0.85,transform shape]
\draw (-1,0) node[left]{$a$} to[short,o-*] (0,0) -- (0,1.8) to[R,l=$R$] (2.2,1.8) to[R,l=$R$] (4.4,1.8) to[R,l=$R$] (6.6,1.8) -- (6.6,0);
\draw (0,0) -- (0,-1.8) to[R,l_=$R$] (2.2,-1.8) to[R,l_=$R$] (4.4,-1.8) to[R,l_=$R$] (6.6,-1.8) -- (6.6,0);
\draw (6.6,0) to[short,*-o] (7.6,0) node[right]{$b$};
\fill (2.2,1.8) circle(2pt) node[above]{$c$}; \fill (4.4,1.8) circle(2pt) node[above]{$d$};
\fill (2.2,-1.8) circle(2pt) node[below]{$e$}; \fill (4.4,-1.8) circle(2pt) node[below]{$f$};
\draw (2.2,1.8) to[nos,l_=$K_1$] (2.2,0.3) -- (4.4,-1.8);
\draw (4.4,1.8) to[nos,l=$K_2$] (4.4,0.3) -- (2.2,-1.8);
\end{circuitikz}}

\newcommand{\circSdosCap}{%
\begin{circuitikz}[scale=0.9,transform shape]
\draw (0,0) to[battery1,invert,l=$\fem_1$] (0,1.6) to[R,l=$R_1$] (0,3.4) -- (6,3.4);
\draw (3,0) to[battery1,invert,l=$\fem_3$] (3,1.6) to[C,l=$C$] (3,3.4);
\draw (6,0) to[battery1,invert,l_=$\fem_2$] (6,1.6) to[R,l_=$R_2$] (6,3.4);
\draw (0,0) -- (6,0);
\fill (3,3.4) circle(2pt) node[above]{$A$}; \fill (3,0) circle(2pt) node[below]{$B$};
\end{circuitikz}}

\fancyhead[L]{\small Física II -- Solucionario de los simulacros}
\fancyhead[R]{\small Material de tutoría}

\newcommand{\sol}[1]{\par\medskip\noindent{\color{azul}\large\bfseries #1}\par\nopagebreak\smallskip}
\newcommand{\simu}[2]{\newpage\section*{Simulacro #1 -- #2}\addcontentsline{toc}{section}{Simulacro #1 -- #2}}
\newcommand{\clave}[1]{\begin{profe}[Clave para la nota completa]#1\end{profe}}

\begin{document}
\begin{center}
{\LARGE\bfseries Solucionario de los tres simulacros}\\[4pt]
{\large Tercer parcial de Física II -- Corriente eléctrica continua}
\end{center}

\begin{alerta}[Úsalo bien]
No leas una solución antes de haber intentado la pregunta durante al menos 15 minutos. Leer soluciones da una falsa sensación de dominio; lo que fija el aprendizaje es el intento previo (práctica de recuperación). Si te trabas, lee solo la primera línea de la solución y vuelve a intentarlo.
\end{alerta}

\section*{Rúbrica de calificación (25 puntos por pregunta)}
\addcontentsline{toc}{section}{Rúbrica}
\begin{center}
\begin{tabularx}{\linewidth}{@{}lcX@{}}
\toprule
\textbf{Criterio} & \textbf{Puntos} & \textbf{Qué se exige}\\
\midrule
Planteamiento físico & 8 & Ley o principio correcto, notación declarada, figura con sistema de referencia, simetría justificada, sentidos de corriente definidos.\\
Desarrollo matemático & 10 & Pasos completos y justificados; integrales con límites correctos; sistema de Kirchhoff independiente y bien resuelto.\\
Resultado y verificación & 4 & Resultado enmarcado, unidades SI, cifras razonables y una comprobación (dimensiones, caso límite o balance de potencias).\\
Claridad & 3 & Orden, letra legible, conclusiones en palabras.\\
\bottomrule
\end{tabularx}
\end{center}
\textbf{Interpretación}: 90--100: listo para el parcial. 75--89: repasa los errores del registro y pasa al siguiente simulacro. Menos de 75: repite las preguntas falladas dentro de 2 o 3 días antes de avanzar.

\tableofcontents

\simu{1}{confianza}

%--------------------------------------------------------------
\subsection*{Grupo I}
\sol{1.- Corriente, intensidad y densidad de corriente}
\textbf{Corriente eléctrica}: movimiento ordenado de portadores de carga bajo la acción de un campo eléctrico. \textbf{Intensidad}: carga que atraviesa una sección por unidad de tiempo, $I=dq/dt$ ($\mathrm{A}=\mathrm{C/s}$). \textbf{Sentido convencional}: el del movimiento de las cargas positivas, opuesto al de los electrones en un metal.

\textbf{Densidad de corriente}: vector $\vj$ con la dirección de la velocidad de las cargas positivas y módulo $j=dI/dS_\perp$ ($\mathrm{A/m^2}$).

\textbf{Demostración}: sea $d\vec S=\hat n\,dS$ un elemento de superficie que forma un ángulo $\theta$ con $\vj$. Su proyección perpendicular a $\vj$ es $dS_\perp=dS\cos\theta$, de modo que $dI=j\,dS\cos\theta=\vj\cdot d\vec S$. Sumando sobre toda la superficie:
\[
\boxed{I=\int_S\vj\cdot d\vec S}
\]
La intensidad es el flujo de $\vj$. Para $\vj$ uniforme y normal a una sección $A$: $I=jA$.
\clave{Figura con $\vj$, $d\vec S$ y el ángulo $\theta$, y explicar por qué se usa la proyección $dS\cos\theta$.}

\sol{2.- $\vj=nq\vec v_d$ y velocidad de deriva}
\textbf{Deducción}: en un tiempo $dt$ cruzan $dS_\perp$ los portadores contenidos en un cilindro de base $dS_\perp$ y altura $v_d\,dt$: $dN=n\,v_d\,dt\,dS_\perp$. La carga es $dq=q\,dN$, así que $j=\dfrac{dq}{dt\,dS_\perp}=nqv_d$ y, vectorialmente, $\vj=nq\vec v_d$.

\textbf{Cálculo}: $A=\pi(1.00\times10^{-3})^2=3.14\times10^{-6}\uu{m^2}$.
\[
j=\frac{I}{A}=\frac{3.00}{3.14\times10^{-6}}=9.55\times10^{5}\uu{A/m^2},
\]
\[
v_d=\frac{j}{ne}=\frac{9.55\times10^5}{(8.50\times10^{28})(1.602\times10^{-19})}=7.01\times10^{-5}\uu{m/s}.
\]
Tiempo para recorrer $1.00\uu{m}$: $t=1.00/7.01\times10^{-5}=1.43\times10^4\uu{s}\approx3.96\uu{h}$.

\textbf{Interpretación}: la lámpara enciende al instante porque el campo eléctrico se establece en todo el circuito casi a la velocidad de la luz; los electrones de todo el cable empiezan a derivar a la vez.

\sol{3.- Conductividad, resistividad y temperatura}
\textbf{Conductividad} $\sigma$: constante de proporcionalidad entre $\vj$ y $\vE$, $\vj=\sigma\vE$ ($\mathrm{S/m}$). \textbf{Resistividad} $\rho=1/\sigma$ ($\Omega\cdot\mathrm{m}$): propiedad del material, no de la forma del conductor.

\textbf{Temperatura}: al calentarse un metal aumenta la amplitud de vibración de los iones de la red, los electrones chocan con más frecuencia, disminuye el tiempo medio entre choques $\tau$ y, como $\sigma=ne^2\tau/m$, aumenta $\rho$. En un intervalo moderado, $\rho=\rho_0[1+\alpha(T-T_0)]$.

\textbf{Cálculo} (despreciando la dilatación):
\[
R_{80}=10.0\,[1+3.90\times10^{-3}(60.0)]=\boxed{12.3\ohm},\qquad
T=20.0+\frac{12.0/10.0-1}{3.90\times10^{-3}}=\boxed{71.3\,^\circ\mathrm{C}}.
\]

\sol{4.- De la ley de Ohm diferencial a la integral}
\textbf{Ley de Ohm diferencial}: en cada punto de un conductor isótropo, $\vj=\sigma\vE$.

\textbf{Deducción}: en un cilindro homogéneo de longitud $L$ y sección $A$ con sus extremos a potenciales $\varphi_1>\varphi_2$, por simetría el campo es uniforme y axial: $E=(\varphi_1-\varphi_2)/L=U/L$. Entonces $j=\sigma U/L$ es uniforme y
\[
I=\int_A\vj\cdot d\vec S=jA=\frac{\sigma A}{L}U\ \Longrightarrow\ \boxed{I=\frac UR},\qquad\boxed{R=\frac{L}{\sigma A}=\rho\frac LA}.
\]
\textbf{Verificación de unidades}: $[\rho L/A]=\Omega\cdot\mathrm m\cdot\mathrm m/\mathrm{m^2}=\Omega$.

\sol{5.- Fem, diferencia de potencial y caída de tensión; circuito cerrado}
\textbf{Fem}: trabajo que realizan las fuerzas extrañas (no electrostáticas) para llevar la unidad de carga positiva a través de la fuente, del borne $-$ al $+$: $\fem=A^*/q=\oint\vE^*\cdot d\vec l$.

\textbf{Diferencias}: la \emph{diferencia de potencial} es el trabajo del campo electrostático por unidad de carga; la \emph{fem}, el de las fuerzas extrañas; la \emph{caída de tensión} $U_{12}=IR_{12}$, el trabajo total: $U_{12}=(\varphi_1-\varphi_2)+\fem_{12}$. En un tramo sin fuentes, caída de tensión y diferencia de potencial coinciden.

\textbf{Deducción}: aplicando la ley de Ohm para un tramo no homogéneo a todo el circuito (desde un punto hasta él mismo, $\varphi_1=\varphi_2$), con resistencia total $R+r$: $I(R+r)=\fem$, de donde
\[
\boxed{I=\frac{\fem}{R+r}},\qquad\boxed{U_{\text{bornes}}=IR=\fem-Ir}.
\]
Casos límite: circuito abierto ($I=0$), $U=\fem$; cortocircuito ($R=0$), $I_{cc}=\fem/r$.

%--------------------------------------------------------------
\subsection*{Grupo II}
\sol{6.- Alambre de nicrom}
$A=\pi(0.250\times10^{-3})^2=1.96\times10^{-7}\uu{m^2}$.
\[
R=\rho\frac LA=\frac{(1.10\times10^{-6})(2.00)}{1.96\times10^{-7}}=11.2\ohm,\quad I=\frac{12.0}{11.2}=1.07\uu A,
\]
\[
j=\frac IA=5.46\times10^6\uu{A/m^2},\qquad E=\frac UL=\frac{12.0}{2.00}=6.00\uu{V/m}.
\]
\textbf{Verificación}: $\rho j=(1.10\times10^{-6})(5.46\times10^6)=6.00\uu{V/m}=E$. \checkmark

\sol{7.- Batería real}
\[
I=\frac{\fem}{R+r}=\frac{12.0}{6.00}=2.00\uu A,\qquad U=\fem-Ir=12.0-1.00=11.0\uu V.
\]
$P_\fem=\fem I=24.0\uu W$; $P_r=I^2r=2.00\uu W$; $P_R=I^2R=22.0\uu W$. Balance: $24.0=2.00+22.0$ \checkmark. Rendimiento: $\eta=R/(R+r)=U/\fem=91.7\%$.

\sol{8.- Asociaciones}
(a) $R_s=2.0+3.0+6.0=\boxed{11.0\ohm}$. \quad (b) $\dfrac1{R_p}=\dfrac12+\dfrac13+\dfrac16=1\Rightarrow\boxed{1.00\ohm}$. \quad (c) $R_{23}=\dfrac{3\cdot6}{9}=2.0\ohm$, $R=2.0+2.0=\boxed{4.00\ohm}$.

Con $12.0\uu V$ en (c): $I_1=12.0/4.00=3.00\uu A$; la tensión sobre el paralelo es $3.00\cdot2.0=6.00\uu V$; $I_2=6.00/3.0=2.00\uu A$; $I_3=6.00/6.0=1.00\uu A$. Nodo: $2.00+1.00=3.00$ \checkmark.

\sol{9.- Serie y paralelo desde Kirchhoff}
\textbf{Serie}: entre resistencias consecutivas no hay derivaciones, así que por la regla de nodos la corriente $I$ es la misma en todas. Aplicando la regla de mallas al lazo formado con la fuente: $U=\sum IR_k=I\sum R_k$. El equivalente debe cumplir $U=IR_s$, luego $R_s=\sum R_k$.

\textbf{Paralelo}: todas comparten los mismos dos nodos; la regla de mallas aplicada al lazo formado por dos resistencias da $I_iR_i=I_jR_j=U$. La regla de nodos da $I=\sum I_k=U\sum 1/R_k$, y como $I=U/R_p$: $1/R_p=\sum1/R_k$.

\textbf{$R_p<R_{\min}$}: todos los términos son positivos, así que $\dfrac1{R_p}=\dfrac1{R_{\min}}+(\text{términos}>0)>\dfrac1{R_{\min}}$, es decir, $R_p<R_{\min}$. Físicamente, cada rama nueva añade un camino para la corriente.

\sol{10.- Calentador}
$R=U^2/P=120^2/1000=14.4\ohm$; $I=P/U=8.33\uu A$; $E=Pt=1.00\uu{kW}\times0.500\uu h=0.500\uu{kWh}=1.80\times10^6\uu J$.

A $110\uu V$: $P'=110^2/14.4=840\uu W$. \textbf{Suposición}: $R$ constante. En realidad el elemento se calienta un poco menos y su $R$ baja ligeramente, de modo que $P'$ real es algo mayor que $840\uu W$.

%--------------------------------------------------------------
\subsection*{Grupo III}
\sol{11.- Reglas de Kirchhoff}
\textbf{1.ª regla (nodos)}: $\sum_kI_k=0$ en cada nodo. Se fundamenta en la conservación de la carga en régimen estacionario: aplicando $\oint\vj\cdot d\vec S=-dq/dt=0$ a una superficie que encierra el nodo, solo la atraviesan los conductores, así que la suma algebraica de sus corrientes es cero.

\textbf{2.ª regla (mallas)}: $\sum_kI_kR_k=\sum_k\fem_k$ en cada contorno cerrado. Se fundamenta en la ley de Ohm para un tramo no homogéneo, $I_kR_k=\varphi_k-\varphi_{k+1}+\fem_k$, sumada a lo largo del contorno: las diferencias de potencial se cancelan porque el campo electrostático es conservativo ($\oint\vE\cdot d\vec l=0$). Expresa la conservación de la energía.

\textbf{Procedimiento}: (1) asignar una corriente con sentido arbitrario a cada una de las $m$ ramas; (2) escribir $n-1$ ecuaciones de nodos; (3) escribir $m-n+1$ ecuaciones de malla en mallas independientes, cada una con al menos una rama nueva (por ejemplo, las «ventanas» de un circuito plano); (4) convenio de signos: $+IR$ si se recorre a favor de la corriente, $+\fem$ si se pasa de $-$ a $+$; (5) resolver; una corriente negativa indica que su sentido real es el opuesto; (6) verificar con el balance de potencias.

\sol{12.- Circuito de dos mallas}
\begin{center}\circSunoQdoce\end{center}
\textbf{Sentidos elegidos}: $I_1$ hacia arriba por $\fem_1$ (de $B$ hacia $A$), $I_3$ hacia arriba por $\fem_2$, $I_2$ hacia abajo por $R_2$ (de $A$ a $B$).

\textbf{Nodo $A$}: $I_1+I_3=I_2$.

\textbf{Malla izquierda} (horaria: sube por $\fem_1$ y baja por $R_2$): $\fem_1=I_1R_1+I_2R_2\Rightarrow12=2I_1+6I_2$.

\textbf{Malla derecha} (sube por $\fem_2$, pasa por $R_3$ hacia $A$ y baja por $R_2$): $\fem_2=I_3R_3+I_2R_2\Rightarrow6=3I_3+6I_2$.

\textbf{Solución}: $I_1=6-3I_2$ e $I_3=2-2I_2$. Sustituyendo en el nodo: $8-5I_2=I_2\Rightarrow$
\[
\boxed{I_2=\tfrac43\uu A=1.33\uu A},\qquad\boxed{I_1=2.00\uu A},\qquad\boxed{I_3=-\tfrac23\uu A=-0.667\uu A}.
\]
$I_3<0$: por $\fem_2$ la corriente baja, es decir, \emph{entra por su borne positivo}: la fuente 2 se está cargando. $\varphi_A-\varphi_B=I_2R_2=\boxed{8.00\uu V}$.

\textbf{Balance de potencias}: $\fem_1$ entrega $12(2)=24.0\uu W$; $\fem_2$ \emph{absorbe} $6(2/3)=4.00\uu W$. Joule: $R_1$: $4(2)=8.00$; $R_2$: $(16/9)(6)=10.67$; $R_3$: $(4/9)(3)=1.33$; total $20.0\uu W$. Se cumple $24.0=20.0+4.00$ \checkmark.

\textbf{Verificación con potenciales de nodo} (Millman): $\varphi_A-\varphi_B=\dfrac{12/2+6/3}{1/2+1/6+1/3}=8.00\uu V$ \checkmark.

\sol{13.- Puente de Wheatstone}
Con $I_G=0$, la corriente $I_1$ pasa por $R_1$ y $R_2$, e $I_3$ por $R_3$ y $R_x$. Equilibrio: $\varphi_c=\varphi_d$, luego $I_1R_1=I_3R_3$ (caídas desde $a$) e $I_1R_2=I_3R_x$ (caídas hasta $b$). Dividiendo:
\[
\frac{R_1}{R_2}=\frac{R_3}{R_x}\ \Rightarrow\ R_x=\frac{R_2R_3}{R_1}=\frac{150\cdot200}{100}=\boxed{300\ohm}.
\]
La condición no depende de la fem ni de la resistencia del galvanómetro.

\sol{14.- Amperímetro y voltímetro}
(a) \textbf{Shunt en paralelo} $R_s$: misma tensión en ambos, $I_gR_g=(I-I_g)R_s$:
\[
R_s=\frac{I_gR_g}{I-I_g}=\frac{(1.00\times10^{-3})(50.0)}{2.00-0.001}=\boxed{0.0250\ohm}\ \text{en paralelo}.
\]
(b) \textbf{Resistencia en serie} $R_m$: misma corriente, $U=I_g(R_g+R_m)$:
\[
R_m=\frac{U}{I_g}-R_g=10\,000-50.0=\boxed{9950\ohm}\ \text{en serie}.
\]
Un amperímetro ideal tiene $R\to0$ y un voltímetro ideal, $R\to\infty$, para no perturbar el circuito.

\sol{15.- Puente de resistencias iguales}
\begin{center}\circPuente{$R$}{$R$}{$R$}{$R$}{$R$}\end{center}
\textbf{Simetría}: la reflexión respecto del eje $a$--$b$ intercambia $c$ y $d$ sin alterar $a$ ni $b$, así que $\varphi_c=\varphi_d$ (también es un puente equilibrado: $R/R=R/R$). Por la resistencia $c$--$d$ no circula corriente y se puede quitar:
\[
R_{eq}=(R+R)\,\|\,(R+R)=\boxed{R}.
\]
Si se sustituye por un cable, el resultado \textbf{no cambia}: unir dos puntos que ya están al mismo potencial no modifica ninguna corriente. Comprobación: $(R\|R)+(R\|R)=R/2+R/2=R$ \checkmark.

%--------------------------------------------------------------
\subsection*{Grupo IV}
\sol{16.- Ley de Joule--Lenz}
\textbf{Integral}: al pasar $dq=I\,dt$ por un conductor con diferencia de potencial $U$, el campo realiza el trabajo $dA=U\,dq=UI\,dt$. Si el conductor está fijo y no hay transformaciones químicas, ese trabajo se transforma íntegramente en calor: $dQ=UI\,dt=I^2R\,dt$, de donde $Q=\int I^2R\,dt$, y con $I$ constante, $Q=I^2Rt$.

\textbf{Diferencial}: en un cilindro elemental de longitud $dl$ y sección $dS$ orientado según $\vj$: $dR=\rho\,dl/dS$, $dI=j\,dS$, así que $dP=dI^2\,dR=\rho j^2\,dV$. Por lo tanto
\[
\boxed{w=\frac{dP}{dV}=\rho j^2=\vj\cdot\vE=\sigma E^2}.
\]
\textbf{Significado}: $w$ es la potencia térmica por unidad de volumen ($\mathrm{W/m^3}$), es decir, el calor que se libera por segundo en cada metro cúbico del conductor. Es una ley local, válida aunque el conductor no sea homogéneo.

\sol{17.- Máxima transferencia de potencia}
$P(R)=I^2R=\dfrac{\fem^2R}{(R+r)^2}$. Derivando: $\dfrac{dP}{dR}=\fem^2\dfrac{(R+r)^2-2R(R+r)}{(R+r)^4}=\fem^2\dfrac{r-R}{(R+r)^3}=0\Rightarrow\boxed{R=r}$. Es un máximo porque $dP/dR>0$ para $R<r$ y $dP/dR<0$ para $R>r$.
\[
P_{\max}=\frac{\fem^2}{4r},\qquad\eta=\frac{R}{R+r}=50\%.
\]
Con $\fem=9.0\uu V$ y $r=1.5\ohm$: $R=1.5\ohm$, $I=3.0\uu A$, $P_{\max}=81/6=\boxed{13.5\uu W}$.

\sol{18.- Bombillas en serie}
$R_{60}=120^2/60=240\ohm$, $R_{100}=120^2/100=144\ohm$. En serie: $I=120/384=0.3125\uu A$.
\[
P_{60}=I^2(240)=\boxed{23.4\uu W},\qquad P_{100}=I^2(144)=\boxed{14.1\uu W}.
\]
\textbf{Brilla más la de $60\uu W$}: en serie la corriente es común y disipa más la de mayor resistencia.

\textbf{Suposición}: resistencias constantes. No es realista: el filamento de tungsteno frío tiene una resistencia varias veces menor que caliente, y aquí trabajan por debajo de su temperatura nominal. Las potencias reales difieren de estas, pero la conclusión cualitativa se mantiene en la práctica.

\sol{19.- Capacitor en régimen estacionario}
\begin{center}\circSunoQdiecinueve\end{center}
En régimen estacionario no pasa corriente por la rama del capacitor:
\[
I=\frac{\fem}{r+R_1+R_2}=\frac{12.0}{6.0}=2.00\uu A,\qquad U_C=IR_2=4.00\uu V,
\]
\[
Q=CU_C=40.0\uu{\mu C},\quad W=\tfrac12CU_C^2=80.0\uu{\mu J}.
\]
\textbf{Sin $R_2$}: no hay ningún camino cerrado para la corriente continua, así que $I=0$, no hay caída en $r$ ni en $R_1$ y $U_C=\fem=12.0\uu V$. Entonces $Q'=\boxed{120\uu{\mu C}}$.

\sol{20.- Cuadrado con diagonal}
\begin{center}\circCuadrado\end{center}
\textbf{Entre $A$ y $C$}: tres caminos en paralelo: $A$--$B$--$C$ ($2R$), $A$--$D$--$C$ ($2R$) y la diagonal ($R$):
\[
\frac1{R_{AC}}=\frac1{2R}+\frac1{2R}+\frac1R=\frac2R\ \Rightarrow\ \boxed{R_{AC}=R/2}.
\]
\textbf{Entre $B$ y $D$}: la reflexión respecto de la recta $BD$ intercambia $A$ y $C$ sin mover $B$ ni $D$, así que $\varphi_A=\varphi_C$ y por la diagonal no circula corriente:
\[
R_{BD}=(2R)\,\|\,(2R)=\boxed{R}.
\]

\simu{2}{nivel del profesor}

%--------------------------------------------------------------
\subsection*{Grupo I}
\sol{1.- Ecuación de continuidad}
\textbf{Principio}: la carga eléctrica se conserva. Sea $q$ la carga dentro de una superficie cerrada fija $S$. Si por $S$ sale una corriente neta $I_{\text{sal}}=\oint_S\vj\cdot d\vec S$, la carga interior disminuye al mismo ritmo:
\[
\boxed{\oint_S\vj\cdot d\vec S=-\frac{dq}{dt}}.
\]
\textbf{Forma diferencial}: con $q=\int_V\rho_q\,dV$ y el teorema de Gauss--Ostrogradski, $\int_V\big(\nabla\cdot\vj+\partial\rho_q/\partial t\big)dV=0$ para todo volumen $V$; como el volumen es arbitrario, el integrando se anula:
\[
\boxed{\nabla\cdot\vj+\frac{\partial\rho_q}{\partial t}=0}.
\]
\textbf{Corriente estacionaria}: la distribución de carga no cambia, $\partial\rho_q/\partial t=0$, así que $\nabla\cdot\vj=0$ y $\oint\vj\cdot d\vec S=0$: las líneas de corriente no tienen principio ni fin (son cerradas).

\textbf{Primera regla de Kirchhoff}: se toma $S$ rodeando un nodo. Solo la cortan los conductores $k$, cada uno con flujo $\pm I_k$, así que $\oint\vj\cdot d\vec S=\sum_k(\pm I_k)=0$.

\sol{2.- Modelo clásico y ley de Ohm diferencial}
\textbf{Hipótesis}: (i) los electrones de conducción forman un gas de partículas libres; (ii) en cada choque con un ion pierden la velocidad dirigida (la velocidad después del choque tiene dirección aleatoria); (iii) el tiempo medio entre choques es $\tau$, independiente de $E$ para campos moderados.

\textbf{Deducción}: entre choques, $m\,d\vec u/dt=-e\vE$, luego $\vec u(t)=-\dfrac{e\vE}{m}t$, con $t$ el tiempo desde el último choque. Promediando sobre el conjunto de electrones, el tiempo medio transcurrido desde el último choque es $\tau$:
\[
\vec v_d=\langle\vec u\rangle=-\frac{e\tau}{m}\vE\ \Longrightarrow\ \vj=-ne\vec v_d=\frac{ne^2\tau}{m}\vE\equiv\sigma\vE,\qquad\boxed{\sigma=\frac{ne^2\tau}{m}}.
\]
(Si se supone que todos los electrones tienen el mismo tiempo libre $\tau_0$, la velocidad media es la mitad de la máxima y resulta $\sigma=ne^2\tau_0/2m$. La ley es la misma, solo cambia el factor numérico; usa el convenio de tu profesor.)

\textbf{Temperatura}: $n$ es prácticamente constante en un metal. Al subir $T$ aumentan las vibraciones de la red, los choques son más frecuentes, $\tau$ disminuye, $\sigma$ disminuye y $\rho=1/\sigma$ aumenta.

\sol{3.- Fuerzas extrañas y ley de Ohm generalizada}
\textbf{Fuerzas extrañas}: fuerzas de origen no electrostático (químicas, magnéticas, térmicas, etc.) que actúan sobre los portadores dentro de una fuente y los llevan en contra del campo electrostático. Se describen con el campo $\vE^*=\vec F^*/q$.

\textbf{Fem}: $\fem_{12}=\dfrac{A^*_{12}}{q}=\displaystyle\int_1^2\vE^*\cdot d\vec l$.

\textbf{Ley de Ohm generalizada (forma diferencial)}: la fuerza total por unidad de carga es $\vE+\vE^*$, así que $\vj=\sigma(\vE+\vE^*)$.

\textbf{Tramo no homogéneo}: despejamos $\vE+\vE^*=\rho\vj$ e integramos a lo largo de un conductor delgado de sección $A$ desde 1 hasta 2, con $\vj\parallel d\vec l$ y $j=I/A$:
\[
\int_1^2\vE\cdot d\vec l+\int_1^2\vE^*\cdot d\vec l=I\int_1^2\frac{\rho\,dl}{A}
\ \Longrightarrow\ (\varphi_1-\varphi_2)+\fem_{12}=IR_{12}.
\]
\[
\boxed{IR_{12}=\varphi_1-\varphi_2+\fem_{12}}
\]

\sol{4.- Imposibilidad de una corriente estacionaria solo electrostática}
Supongamos que en un circuito cerrado $L$ solo actúa el campo electrostático. Entonces $\vj=\sigma\vE$. Recorremos $L$ siguiendo una línea de corriente, de modo que $\vj\parallel d\vec l$ y $\vj\cdot d\vec l=j\,dl>0$:
\[
\oint_L\vE\cdot d\vec l=\oint_L\frac{\vj\cdot d\vec l}{\sigma}=\oint_L\rho j\,dl>0.
\]
Pero el campo electrostático es conservativo: $\oint_L\vE\cdot d\vec l=0$ (segundo teorema electrostático fundamental). La contradicción demuestra que $j=0$: \textbf{sin fuerzas extrañas no hay corriente estacionaria}.

\textbf{Interpretación energética}: en el circuito se disipa continuamente la potencia de Joule $I^2R$, pero el campo electrostático no entrega energía neta en un ciclo cerrado. La fuente transforma otra forma de energía en energía eléctrica a razón de $P=\fem I$. Ejemplos: pila o batería (energía química, fuerzas de origen químico), generador (fuerza de Lorentz sobre las cargas de un conductor en movimiento, inducción), termopar (difusión térmica de portadores), celda solar (energía de los fotones absorbidos).

\sol{5.- Joule--Lenz diferencial y su integración}
\textbf{Deducción}: el campo realiza sobre los portadores contenidos en $dV$ la potencia $dP=\sum\vec F\cdot\vec v=(nq\,dV)\vE\cdot\vec v_d=\vj\cdot\vE\,dV$. En régimen estacionario la energía cinética media no aumenta, así que todo ese trabajo se cede a la red como calor:
\[
\boxed{w=\vj\cdot\vE=\sigma E^2=\rho j^2}.
\]
\textbf{Cilindro homogéneo} (longitud $L$, sección $A$, tensión $U$): $E=U/L$ uniforme, así que
\[
P=\int_Vw\,dV=\sigma\frac{U^2}{L^2}AL=\frac{\sigma A}{L}U^2=\frac{U^2}{R}=I^2R.
\]

%--------------------------------------------------------------
\subsection*{Grupo II}
\sol{6.- Cono truncado}
Eje $x$ desde la base menor; el radio crece linealmente: $r(x)=a+\dfrac{b-a}{L}x$, con $dr=\dfrac{b-a}{L}dx$. Las rebanadas de espesor $dx$ están en serie (aproximación: superficies equipotenciales casi planas, válida si el cono no es muy abierto):
\[
R=\int_0^L\frac{\rho\,dx}{\pi r^2}=\frac{\rho L}{\pi(b-a)}\int_a^b\frac{dr}{r^2}=\frac{\rho L}{\pi(b-a)}\Big(\frac1a-\frac1b\Big)=\boxed{\frac{\rho L}{\pi ab}}.
\]
\textbf{Comprobación}: si $b=a$, $R=\rho L/(\pi a^2)$ \checkmark.

\textbf{Campo}: $E(x)=\rho j=\dfrac{\rho I}{\pi\big(a+\frac{b-a}{L}x\big)^2}$. \textbf{Calor por unidad de longitud}: $\dfrac{dP}{dx}=I^2\dfrac{dR}{dx}=\dfrac{\rho I^2}{\pi r(x)^2}$, máximo en $x=0$ (base menor): el extremo delgado es el que más se calienta.

\sol{7.- Esferas concéntricas}
Corriente radial; cascarones de espesor $dr$ y área $4\pi r^2$ en serie:
\[
R=\int_{R_1}^{R_2}\frac{\rho\,dr}{4\pi r^2}=\boxed{\frac{\rho}{4\pi}\Big(\frac1{R_1}-\frac1{R_2}\Big)=\frac{\rho(R_2-R_1)}{4\pi R_1R_2}},\qquad I=\frac UR=\frac{4\pi R_1R_2\,U}{\rho(R_2-R_1)}.
\]
Del valor de $I$, por continuidad (el flujo de $\vj$ es igual a $I$ en cada esfera de radio $r$):
\[
\vj(r)=\frac{I}{4\pi r^2}\hat r=\frac{R_1R_2U}{\rho(R_2-R_1)}\frac{\hat r}{r^2},\qquad\vE(r)=\rho\vj=\frac{R_1R_2U}{(R_2-R_1)}\frac{\hat r}{r^2}.
\]
\textbf{Comprobación}: $\int_{R_1}^{R_2}E\,dr=U$ \checkmark. El campo es el mismo que en el capacitor esférico electrostático con tensión $U$.

\sol{8.- Fuga en un cable coaxial y $RC=\varepsilon\rho$}
Cascarones cilíndricos de área $2\pi rH$ en serie:
\[
R=\int_{R_1}^{R_2}\frac{\rho\,dr}{2\pi rH}=\boxed{\frac{\rho}{2\pi H}\ln\frac{R_2}{R_1}}.
\]
Con la capacitancia del capacitor cilíndrico (parcial 2), $C=\dfrac{2\pi\varepsilon H}{\ln(R_2/R_1)}$:
\[
RC=\frac{\rho\ln(R_2/R_1)}{2\pi H}\cdot\frac{2\pi\varepsilon H}{\ln(R_2/R_1)}=\boxed{\varepsilon\rho}.
\]
Con tensión $U$: $I=U/R$, $\vj(r)=\dfrac{I}{2\pi rH}\hat r=\dfrac{U}{\rho\,r\ln(R_2/R_1)}\hat r$ y $\vE(r)=\rho\vj=\dfrac{U}{r\ln(R_2/R_1)}\hat r$.

\sol{9.- Resistividad variable y carga en el interior}
\[
R=\int_0^L\frac{\rho_0(1+x/L)}{A}dx=\frac{\rho_0}{A}\Big(L+\frac L2\Big)=\boxed{\frac{3\rho_0L}{2A}}.
\]
Con $j=I/A$ uniforme (continuidad, sección constante), $E(x)=\rho(x)j=\dfrac{\rho_0I}{A}\Big(1+\dfrac xL\Big)$: el campo \emph{no} es uniforme.

\textbf{Gauss}: si se toma como superficie gaussiana un cilindro de sección $A$ entre $x$ y $x+dx$, el flujo neto es $[E(x+dx)-E(x)]A$, así que $\dfrac{dE}{dx}A\,dx=\dfrac{dq}{\varepsilon_0}$. Por lo tanto
\[
\rho_q=\varepsilon_0\frac{dE}{dx}=\boxed{\frac{\varepsilon_0\rho_0I}{AL}}\neq0\quad(\text{uniforme, positiva}).
\]
Carga total en el alambre: $\rho_qAL=\varepsilon_0\rho_0I$. Coincide con la fórmula general $\rho_q=\varepsilon_0\vj\cdot\nabla\rho$. \textbf{Conclusión}: en un conductor no homogéneo con corriente aparece carga volumétrica.

\sol{10.- Electrodo semiesférico de puesta a tierra}
La corriente sale radialmente por semiesferas de área $2\pi r^2$:
\[
R=\int_a^\infty\frac{\rho\,dr}{2\pi r^2}=\boxed{\frac{\rho}{2\pi a}}.
\]
$j(r)=\dfrac{I}{2\pi r^2}$, $E(r)=\dfrac{\rho I}{2\pi r^2}$, $\varphi(r)=\displaystyle\int_r^\infty E\,dr=\dfrac{\rho I}{2\pi r}$ (con $\varphi(\infty)=0$).
\[
U_{\text{paso}}=\varphi(r)-\varphi(r+d)=\frac{\rho I}{2\pi}\Big(\frac1r-\frac1{r+d}\Big)=\boxed{\frac{\rho Id}{2\pi r(r+d)}}.
\]
\textbf{Numérico}: $R=\dfrac{100}{2\pi(0.50)}=31.8\ohm$; potencial del electrodo, $\varphi(a)=IR=3.18\uu{kV}$;
\[
U_{\text{paso}}=\frac{(100)(100)(0.80)}{2\pi(2.0)(2.8)}=\boxed{227\uu V},
\]
peligroso para una persona. (Modelo de suelo homogéneo; los suelos reales son estratificados.)

%--------------------------------------------------------------
\subsection*{Grupo III}
\sol{11.- Dos fuentes en paralelo: expresión general}
Sentidos: $I_1$ e $I_2$ suben por las fuentes hacia $A$; $I_3$ baja por $R_3$.

Nodo $A$: $I_1+I_2=I_3$. Malla 1 ($\fem_1$, $R_1$, $R_3$): $\fem_1=I_1R_1+I_3R_3$. Malla 2 ($\fem_2$, $R_2$, $R_3$): $\fem_2=I_2R_2+I_3R_3$.

Al resolver (por ejemplo, por Cramer) con $D=R_1R_2+R_1R_3+R_2R_3$:
\[
\boxed{I_3=\frac{\fem_1R_2+\fem_2R_1}{D}},\quad\boxed{I_1=\frac{\fem_1(R_2+R_3)-\fem_2R_3}{D}},\quad\boxed{I_2=\frac{\fem_2(R_1+R_3)-\fem_1R_3}{D}}.
\]
Verificación: $I_1+I_2=(\fem_1R_2+\fem_2R_1)/D=I_3$ \checkmark.

\textbf{$I_2=0$} $\iff\fem_2(R_1+R_3)=\fem_1R_3\iff\boxed{\fem_2=\dfrac{\fem_1R_3}{R_1+R_3}}$. Interpretación: la fem 2 iguala exactamente la tensión que $\fem_1$ establece sobre $R_3$ cuando trabaja sola.

\sol{12.- Tres ramas con tres fuentes}
\begin{center}\circSdosQdoce\end{center}
Sentidos: las tres corrientes \emph{hacia arriba} (hacia $A$): $I_1$, $I_2$, $I_3$. Atención: $\fem_3$ tiene el borne positivo hacia $B$.

Nodo $A$: $I_1+I_2+I_3=0$.

Malla izquierda (sube por la rama 1 y baja por la rama 2): $\fem_1-\fem_2=I_1R_1-I_2R_2\Rightarrow4=2I_1-4I_2$.

Malla derecha (sube por la rama 2 y baja por la rama 3; al bajar por $\fem_3$ se pasa de $-$ a $+$): $\fem_2+\fem_3=I_2R_2-I_3R_3\Rightarrow12=4I_2-4I_3$.

Solución: $I_1=2+2I_2$, $I_3=I_2-3$; en el nodo, $4I_2-1=0$:
\[
\boxed{I_2=0.25\uu A},\qquad\boxed{I_1=2.50\uu A},\qquad\boxed{I_3=-2.75\uu A}\ (\text{baja }2.75\uu A).
\]
$\varphi_A-\varphi_B=\fem_1-I_1R_1=12-5=\boxed{7.0\uu V}$ (también $8-0.25\cdot4=7$ y $-4+2.75\cdot4=7$ \checkmark).

\textbf{¿Quién entrega energía?} $\fem_1$: la corriente sale por su $+$, entrega $30.0\uu W$. $\fem_2$: sale por su $+$, entrega $2.0\uu W$. $\fem_3$: la corriente baja y sale por su $+$ (abajo), entrega $4(2.75)=11.0\uu W$. \textbf{Las tres entregan}; ninguna recibe.

\textbf{Balance}: $30+2+11=43\uu W$; Joule: $2.5^2(2)+0.25^2(4)+2.75^2(4)=12.5+0.25+30.25=43.0\uu W$ \checkmark.

\sol{13.- Escalera infinita de resistencias iguales}
\begin{center}\circEscalera{$R$}{$R$}\end{center}
Si se quita la primera celda (la $R$ en serie y la primera $R$ en paralelo), lo que queda a la derecha es otra escalera infinita idéntica, de resistencia $R_\infty$. Por lo tanto la red equivale a $R$ en serie con el paralelo de $R$ y $R_\infty$:
\[
R_\infty=R+\frac{R\,R_\infty}{R+R_\infty}\ \Rightarrow\ R_\infty^2-RR_\infty-R^2=0\ \Rightarrow\ \boxed{R_\infty=\frac{1+\sqrt5}{2}R\approx1.618R}
\]
(se descarta la raíz negativa, sin sentido físico). Es el número áureo.

\sol{14.- Cubo entre vértices opuestos}
Entra $I$ por el vértice $A$ y sale por el opuesto $B$. \textbf{Simetría}: las tres aristas que salen de $A$ son equivalentes (una rotación de $120^\circ$ en torno a la diagonal $AB$ las intercambia), así que cada una lleva $I/3$. Los tres vértices vecinos de $A$ están al mismo potencial; de cada uno salen dos aristas equivalentes con $I/6$ cada una (6 aristas intermedias). Las tres aristas que llegan a $B$ llevan $I/3$.

Por el camino $A\to$ vecino $\to$ segundo vecino $\to B$:
\[
U=\frac I3R+\frac I6R+\frac I3R=\frac56IR\ \Rightarrow\ \boxed{R_{eq}=\frac56R}.
\]

\sol{15.- Transformación triángulo--estrella y puente asimétrico}
\textbf{Deducción}: el triángulo ($R_{ab},R_{bc},R_{ca}$) y la estrella ($R_a,R_b,R_c$) son equivalentes si la resistencia entre cada par de terminales, con el tercero libre, es la misma. Con $\Sigma=R_{ab}+R_{bc}+R_{ca}$:
\[
R_a+R_b=\frac{R_{ab}(R_{bc}+R_{ca})}{\Sigma},\quad R_b+R_c=\frac{R_{bc}(R_{ab}+R_{ca})}{\Sigma},\quad R_c+R_a=\frac{R_{ca}(R_{ab}+R_{bc})}{\Sigma}.
\]
Sumando la 1.ª y la 3.ª y restando la 2.ª: $2R_a=\dfrac{2R_{ab}R_{ca}}{\Sigma}$, luego $\boxed{R_a=\dfrac{R_{ab}R_{ca}}{\Sigma}}$, y análogamente $R_b=\dfrac{R_{ab}R_{bc}}{\Sigma}$ y $R_c=\dfrac{R_{bc}R_{ca}}{\Sigma}$.

\textbf{Aplicación}: el triángulo $a$--$c$--$d$ ($R_{ac}=1$, $R_{cd}=1$, $R_{da}=2$, $\Sigma=4\ohm$) se transforma en una estrella de centro $O$:
\[
R_{aO}=\frac{1\cdot2}{4}=0.5\ohm,\quad R_{cO}=\frac{1\cdot1}{4}=0.25\ohm,\quad R_{dO}=\frac{2\cdot1}{4}=0.5\ohm.
\]
Entonces $R_{ab}=0.5+\big[(0.25+2)\,\|\,(0.5+1)\big]=0.5+\dfrac{2.25\cdot1.5}{3.75}=0.5+0.9=\boxed{1.4\ohm}$.

\textbf{Verificación por nodos} ($\varphi_a=U$, $\varphi_b=0$; por antisimetría, $\varphi_d=U-\varphi_c$): el nodo $c$ da $\varphi_c=4U/7$; la corriente total es $3U/7+2U/7=5U/7$, así que $R=7/5=1.4\ohm$ \checkmark.

%--------------------------------------------------------------
\subsection*{Grupo IV}
\sol{16.- Agrupación óptima de fuentes}
Cada fila equivale a una fuente $(m\fem,\,mr)$; $k$ filas iguales en paralelo equivalen a $(m\fem,\,mr/k)$:
\[
I=\frac{m\fem}{R+mr/k}=\frac{m\fem}{R+m^2r/N}\quad(k=N/m).
\]
\[
\frac{dI}{dm}=\fem\frac{(R+m^2r/N)-m(2mr/N)}{(R+m^2r/N)^2}=0\ \Rightarrow\ R=\frac{m^2r}{N}=\boxed{\frac{mr}{k}=R}.
\]
La corriente es máxima cuando la resistencia interna de la batería equivalente es igual a la externa. En ese caso $m=\sqrt{NR/r}$ e $I_{\max}=\dfrac{\fem}{2}\sqrt{\dfrac{N}{rR}}$. (Si $m$ no resulta entero o no divide a $N$, se comparan los divisores de $N$ más cercanos.)

\sol{17.- Potencia disipada entre esferas concéntricas}
Del problema 7: $j(r)=\dfrac{I}{4\pi r^2}$, así que
\[
w(r)=\rho j^2=\frac{\rho I^2}{16\pi^2r^4}.
\]
Integrando en cascarones de volumen $dV=4\pi r^2dr$:
\[
P=\int_{R_1}^{R_2}\frac{\rho I^2}{16\pi^2r^4}4\pi r^2dr=\frac{\rho I^2}{4\pi}\Big(\frac1{R_1}-\frac1{R_2}\Big)=I^2R=\boxed{\frac{U^2}{R}}.
\]
$w\propto r^{-4}$: el calor se concentra junto a la esfera interior ($r=R_1$).

\sol{18.- Interruptor en un puente}
\begin{center}\circInterruptor\end{center}
\textbf{$K$ abierto}: dos ramas de $3R$ en paralelo: $R_{ab}=\dfrac{3R}{2}$.

\textbf{$K$ cerrado}: $c$ y $d$ forman un solo nodo: $R_{ab}=(R\|2R)+(2R\|R)=\dfrac{2R}{3}+\dfrac{2R}{3}=\dfrac{4R}{3}$.

\[
\Delta R=\frac{4R}{3}-\frac{3R}{2}=\boxed{-\frac R6}\quad(\text{disminuye en }R/6).
\]
\textbf{Corriente por $K$} ($\varphi_a=U$, $\varphi_b=0$): las dos mitades valen $2R/3$, así que $\varphi_c=\varphi_d=U/2$. Llega a $c$ desde $a$: $(U/2)/R=U/2R$; sale de $c$ hacia $b$: $(U/2)/2R=U/4R$. Por la regla de nodos en $c$:
\[
I_K=\frac{U}{2R}-\frac{U}{4R}=\boxed{\frac{U}{4R}}\quad\text{de }c\text{ hacia }d.
\]

\sol{19.- Capacitores, resistencias e interruptor}
\begin{center}\circSdosQdiecinueve\end{center}
Potencial de referencia en el riel inferior ($0$); riel superior a $18\uu V$. Por la rama resistiva circula $I=18/9=2.0\uu A$, de modo que $\varphi_c=18-2(3)=12\uu V$.

\textbf{$K$ abierto}: $C_1$ y $C_2$ están en serie y tienen la misma carga: $C_s=\dfrac{2\cdot4}{6}=\dfrac43\uu{\mu F}$, $Q=C_s\fem=24\uu{\mu C}$. $U_1=12\uu V$, $U_2=6\uu V$, $\varphi_d=6\uu V$.

\textbf{$K$ cerrado}: en régimen estacionario no pasa corriente por $K$, así que la corriente por $R_1$ y $R_2$ sigue siendo $2\uu A$ y $\varphi_d=\varphi_c=12\uu V$:
\[
U_1=6\uu V\Rightarrow Q_1=\boxed{12\uu{\mu C}},\qquad U_2=12\uu V\Rightarrow Q_2=\boxed{48\uu{\mu C}}.
\]
\textbf{Carga por $K$}: el nodo $d$ está formado por la placa inferior de $C_1$ y la superior de $C_2$. Su carga neta pasa de $-24+24=0$ a $-12+48=+36\uu{\mu C}$. Por lo tanto pasan $\boxed{36\uu{\mu C}}$ de $c$ a $d$ a través de $K$.

\sol{20.- Tetraedro}
\begin{center}\circTetraedro\end{center}
Entre los vértices $1$ y $2$: la reflexión que deja fijos $1$ y $2$ intercambia $3$ y $4$, así que $\varphi_3=\varphi_4$ y por la arista $3$--$4$ no circula corriente. Quedan tres caminos en paralelo: la arista directa ($R$), $1$--$3$--$2$ ($2R$) y $1$--$4$--$2$ ($2R$):
\[
\frac1{R_{eq}}=\frac1R+\frac1{2R}+\frac1{2R}=\frac2R\ \Rightarrow\ \boxed{R_{eq}=\frac R2}.
\]
Quitar la arista $3$--$4$ \textbf{no cambia nada}, porque por ella no pasaba corriente. Es un puente de Wheatstone equilibrado.

\simu{3}{superior al profesor}

%--------------------------------------------------------------
\subsection*{Grupo I}
\sol{1.- Ley de Ohm generalizada y balance energético}
Sobre un portador actúan la fuerza electrostática $q\vE$ y la extraña $q\vE^*$. Repitiendo el razonamiento del modelo clásico con la fuerza total, se obtiene $\vj=\sigma(\vE+\vE^*)$.

\textbf{Circuito cerrado}: se despeja $\vE+\vE^*=\rho\vj$ y se integra a lo largo de todo el circuito (conductor delgado, $\vj\parallel d\vec l$, $j=I/A$):
\[
\underbrace{\oint\vE\cdot d\vec l}_{=0}+\underbrace{\oint\vE^*\cdot d\vec l}_{=\fem}=I\oint\frac{\rho\,dl}{A}=I(R+r)\ \Rightarrow\ \boxed{I=\frac{\fem}{R+r}}.
\]
\textbf{Energía}: multiplicando por $I$, $\fem I=I^2R+I^2r$. El término $\fem I$ es la potencia de las fuerzas extrañas (energía no eléctrica que se convierte en eléctrica por segundo). Por la ley de Joule--Lenz, $I^2R$ es el calor que se libera en el circuito externo e $I^2r$ el que se libera dentro de la fuente. El trabajo neto del campo electrostático en un ciclo es nulo: solo transporta la energía desde la fuente hasta la carga.
\[
\eta=\frac{I^2R}{\fem I}=\boxed{\frac{R}{R+r}=\frac{U}{\fem}}.
\]

\sol{2.- Carga en el interior de conductores con corriente}
\textbf{Homogéneo}: régimen estacionario, $\nabla\cdot\vj=0$; con $\sigma$ constante, $\nabla\cdot\vE=\nabla\cdot\vj/\sigma=0$; por Gauss, $\rho_q=\varepsilon_0\nabla\cdot\vE=0$. Las cargas que crean el campo están en la superficie del conductor y en las fuentes.

\textbf{No homogéneo}: $\vE=\rho\vj$, así que $\nabla\cdot\vE=\nabla\cdot(\rho\vj)=\rho\,\nabla\cdot\vj+\vj\cdot\nabla\rho=\vj\cdot\nabla\rho$. Por lo tanto
\[
\boxed{\rho_q=\varepsilon_0\,\vj\cdot\nabla\rho}.
\]
\textbf{Interfaz} (corriente normal, de 1 a 2): por continuidad, $j_1=j_2=j$, así que $E_1=\rho_1j$ y $E_2=\rho_2j$. Con un cilindro gaussiano delgado que atraviesa la interfaz, $\varepsilon_0(E_2-E_1)S=\sigma_sS$:
\[
\boxed{\sigma_s=\varepsilon_0 j(\rho_2-\rho_1)}.
\]
Si los medios tienen permitividades $\varepsilon_1$ y $\varepsilon_2$, la carga \emph{libre} es $\sigma_f=D_2-D_1=j(\varepsilon_2\rho_2-\varepsilon_1\rho_1)$ (ver el problema 8).

\sol{3.- Refracción de las líneas de corriente}
\textbf{Componente normal de $\vj$}: un cilindro aplanado sobre la interfaz y la continuidad estacionaria, $\oint\vj\cdot d\vec S=0$, dan $j_{1n}=j_{2n}$.

\textbf{Componente tangencial de $\vE$}: un rectángulo estrecho que cruza la interfaz y $\oint\vE\cdot d\vec l=0$ dan $E_{1t}=E_{2t}$, es decir, $j_{1t}/\sigma_1=j_{2t}/\sigma_2$.

Con $\alpha_i$ medidos desde la normal, $\tan\alpha_i=j_{it}/j_{in}$:
\[
\frac{\tan\alpha_1}{\tan\alpha_2}=\frac{j_{1t}/j_{1n}}{j_{2t}/j_{2n}}=\frac{j_{1t}}{j_{2t}}=\boxed{\frac{\sigma_1}{\sigma_2}}.
\]
\textbf{Consecuencia}: al pasar de un buen conductor a uno malo ($\sigma_1\gg\sigma_2$), las líneas salen casi perpendiculares a la superficie del buen conductor; por eso los electrodos metálicos se comportan como equipotenciales.

\sol{4.- $RC=\varepsilon\rho$ para cualquier geometría y descarga}
Sea $S$ una superficie cerrada que rodea la armadura positiva, de carga $q$, y que va por dentro del dieléctrico. Corriente de fuga:
\[
I=\oint_S\vj\cdot d\vec S=\frac1\rho\oint_S\vE\cdot d\vec S=\frac{1}{\rho\varepsilon}\oint_S\vD\cdot d\vec S=\frac{q}{\varepsilon\rho}
\]
(dieléctrico homogéneo: $\vD=\varepsilon\vE$; Gauss para $\vD$: flujo igual a la carga libre). Como $R=U/I$ y $C=q/U$:
\[
RC=\frac UI\cdot\frac qU=\frac qI=\boxed{\varepsilon\rho}.
\]
\textbf{Descarga} del capacitor aislado: $\dfrac{dq}{dt}=-I=-\dfrac{q}{\varepsilon\rho}\Rightarrow\boxed{q(t)=q_0e^{-t/\varepsilon\rho}}$. La constante de tiempo $\tau=\varepsilon\rho$ depende solo del material, no de la forma ni del tamaño.

\sol{5.- Joule--Lenz desde el modelo clásico}
\textbf{Modelo de tiempo libre fijo $\tau_0$}: tras cada choque el electrón parte con velocidad dirigida nula y la aumenta linealmente hasta $u_{\max}=eE\tau_0/m$ en el siguiente choque. Velocidad de deriva: $v_d=u_{\max}/2=eE\tau_0/2m$, luego $j=nev_d$ y
\[
\sigma=\frac{ne^2\tau_0}{2m}.
\]
\textbf{Energía cedida por choque}: la energía cinética del movimiento ordenado ganada en el recorrido libre, $\Delta K=\tfrac12mu_{\max}^2=\dfrac{e^2E^2\tau_0^2}{2m}$, se transfiere a la red en el choque.

\textbf{Potencia por unidad de volumen}: hay $n$ electrones por unidad de volumen y cada uno choca $1/\tau_0$ veces por segundo:
\[
w=n\frac{\Delta K}{\tau_0}=\frac{ne^2\tau_0}{2m}E^2=\boxed{\sigma E^2}\quad\text{con la misma }\sigma\text{ de la ley de Ohm.}
\]
(Con la estadística de Drude, $\sigma=ne^2\tau/m$ y $\langle\Delta K\rangle=e^2E^2\langle t^2\rangle/2m=e^2E^2\tau^2/m$; de nuevo $w=\sigma E^2$. El modelo es coherente en ambas versiones.)

%--------------------------------------------------------------
\subsection*{Grupo II}
\sol{6.- Esferas con conductividad variable}
\[
R=\int_{R_1}^{R_2}\frac{dr}{\sigma(r)\,4\pi r^2}=\int_{R_1}^{R_2}\frac{r\,dr}{4\pi\sigma_0R_1r^2}=\boxed{\frac{1}{4\pi\sigma_0R_1}\ln\frac{R_2}{R_1}}.
\]
$j=\dfrac{I}{4\pi r^2}$ y $E=\dfrac j\sigma=\boxed{\dfrac{I}{4\pi\sigma_0R_1\,r}}$, que decrece como $1/r$, no como $1/r^2$: debe haber carga en el medio.

\textbf{Gauss en coordenadas esféricas}: $\nabla\cdot\vE=\dfrac1{r^2}\dfrac{d}{dr}(r^2E)=\dfrac1{r^2}\dfrac{I}{4\pi\sigma_0R_1}$, así que
\[
\rho_q(r)=\boxed{\frac{\varepsilon_0I}{4\pi\sigma_0R_1\,r^2}},\qquad Q_{\text{medio}}=\int_{R_1}^{R_2}\rho_q4\pi r^2dr=\boxed{\frac{\varepsilon_0I(R_2-R_1)}{\sigma_0R_1}}.
\]
\textbf{Comprobación} con $\rho_q=\varepsilon_0\vj\cdot\nabla\rho$: $\rho=r/(\sigma_0R_1)$, $d\rho/dr=1/(\sigma_0R_1)$, así que $\rho_q=\varepsilon_0\dfrac{I}{4\pi r^2}\dfrac{1}{\sigma_0R_1}$ \checkmark.

\sol{7.- Sector de anillo con corriente azimutal}
Por simetría, el potencial depende solo del ángulo: $\varphi=U\theta/\theta_0$, que cumple la ecuación de Laplace en coordenadas polares. El campo es azimutal: $E=\dfrac1r\dfrac{\partial\varphi}{\partial\theta}=\dfrac{U}{\theta_0r}$.

\textbf{Tiras en paralelo}: la tira de radio $r$, ancho $dr$ y espesor $h$ tiene longitud $\theta_0r$ y está sometida a la tensión completa $U$. Su conductancia es $dG=\dfrac{h\,dr}{\rho\,\theta_0r}$:
\[
G=\frac{h}{\rho\theta_0}\int_a^b\frac{dr}{r}=\frac{h}{\rho\theta_0}\ln\frac ba\ \Rightarrow\ \boxed{R=\frac{\rho\,\theta_0}{h\ln(b/a)}}.
\]
\textbf{Comprobación}: si $b=a+\Delta$ con $\Delta\ll a$, $\ln(b/a)\approx\Delta/a$ y $R\approx\rho\,\theta_0a/(h\Delta)=\rho L/A$ \checkmark.

\sol{8.- Capacitor plano con dos capas con pérdidas}
Sentido de la corriente: de la placa 1 (a mayor potencial, en contacto con la capa 1) a la placa 2. Continuidad: la misma $j$ en ambas capas.
\[
U=E_1d_1+E_2d_2=j(\rho_1d_1+\rho_2d_2)\ \Rightarrow\ \boxed{j=\frac{U}{\rho_1d_1+\rho_2d_2}},\quad I=jS,\quad E_1=\rho_1j,\quad E_2=\rho_2j.
\]
\textbf{Cargas libres} (Gauss para $\vD$, con $D_i=\varepsilon_iE_i$):
\[
\sigma_{\text{placa 1}}=+\varepsilon_1\rho_1j,\qquad\sigma_{\text{interfaz}}=D_2-D_1=\boxed{j(\varepsilon_2\rho_2-\varepsilon_1\rho_1)},\qquad\sigma_{\text{placa 2}}=-\varepsilon_2\rho_2j.
\]
La suma es cero (el sistema es neutro) \checkmark. \textbf{Sin carga en la interfaz} $\iff\boxed{\varepsilon_1\rho_1=\varepsilon_2\rho_2}$, es decir, cuando ambos materiales tienen el mismo tiempo de relajación $\varepsilon\rho$.

\sol{9.- Dos electrodos semiesféricos}
Por superposición: el electrodo 1 inyecta $I$ y el 2 la extrae. Para $d\gg a$, cada electrodo es casi equipotencial y el potencial en el suelo es
\[
\varphi=\frac{\rho I}{2\pi}\Big(\frac1{r_1}-\frac1{r_2}\Big),\qquad\varphi_1\approx\frac{\rho I}{2\pi}\Big(\frac1a-\frac1d\Big),\qquad\varphi_2\approx-\varphi_1.
\]
\[
R=\frac{\varphi_1-\varphi_2}{I}\approx\boxed{\frac{\rho}{\pi}\Big(\frac1a-\frac1d\Big)}.
\]
\textbf{Numérico}: $R\approx\dfrac{50}{\pi}\Big(\dfrac1{0.30}-\dfrac1{20}\Big)=15.9\times3.28=\boxed{52\ohm}$.

\textbf{Discusión}: cada electrodo aporta $\rho/(2\pi a)=26.5\ohm$; el término $-1/d$ solo corrige un $1.5\%$. La resistencia de un electrodo hasta un radio $r$ es $\frac{\rho}{2\pi}\big(\frac1a-\frac1r\big)$, que para $r=10a$ ya es el $90\%$ del total: casi toda la resistencia está en las inmediaciones de los electrodos y apenas depende de la distancia entre ellos.

\sol{10.- Descarga de un capacitor esférico con fuga}
Por el problema 4, $\tau=\varepsilon\rho$, $q(t)=q_0e^{-t/\varepsilon\rho}$ e $I(t)=\dfrac{q_0}{\varepsilon\rho}e^{-t/\varepsilon\rho}$.

\textbf{Camino 1 (Joule--Lenz)}: $E=\dfrac{q}{4\pi\varepsilon r^2}$, $w=\dfrac{E^2}{\rho}$:
\[
P(t)=\int_{R_1}^{R_2}\frac{q^2}{16\pi^2\varepsilon^2\rho r^4}4\pi r^2dr=\frac{q^2}{4\pi\varepsilon^2\rho}\Big(\frac1{R_1}-\frac1{R_2}\Big),
\]
\[
Q_J=\int_0^\infty P\,dt=\frac{1}{4\pi\varepsilon^2\rho}\Big(\frac1{R_1}-\frac1{R_2}\Big)q_0^2\frac{\varepsilon\rho}{2}=\boxed{\frac{q_0^2}{8\pi\varepsilon}\Big(\frac1{R_1}-\frac1{R_2}\Big)}.
\]
\textbf{Camino 2 (energía)}: $C=\dfrac{4\pi\varepsilon R_1R_2}{R_2-R_1}$, así que $\dfrac{q_0^2}{2C}=\dfrac{q_0^2(R_2-R_1)}{8\pi\varepsilon R_1R_2}$, el mismo valor \checkmark. Toda la energía almacenada termina como calor en el dieléctrico.

%--------------------------------------------------------------
\subsection*{Grupo III}
\sol{11.- Circuito de cuatro nodos y seis ramas}
\begin{center}\circSdosK\end{center}
\textbf{Corrientes}: $I_1$ ($G\to A$ por $\fem_1$), $I_2$ ($A\to B$), $I_3$ ($A\to C$), $I_4$ ($C\to B$), $I_5$ ($B\to G$), $I_6$ ($G\to C$ por $\fem_2$). Con $n=4$ y $m=6$: 3 ecuaciones de nodo y $6-4+1=3$ de malla.

\textbf{Nodos}: $A$: $I_1=I_2+I_3$; \ $B$: $I_2+I_4=I_5$; \ $C$: $I_3+I_6=I_4$.

\textbf{Mallas} (ventanas):
\begin{itemize}
\item $G\!-\!A\!-\!C\!-\!G$: $24-2I_1-4I_3+2I_6-14=0\Rightarrow2I_1+4I_3-2I_6=10$.
\item $A\!-\!B\!-\!C\!-\!A$: $-2I_2+I_4+4I_3=0$.
\item $C\!-\!B\!-\!G\!-\!C$: $-I_4-2I_5+14-2I_6=0\Rightarrow I_4+2I_5+2I_6=14$.
\end{itemize}
\textbf{Resolución eficiente (potenciales de nodo, equivalente a lo anterior)}: con $\varphi_G=0$,
\[
\begin{aligned}
A:&\ \tfrac{\varphi_A-24}{2}+\tfrac{\varphi_A-\varphi_B}{2}+\tfrac{\varphi_A-\varphi_C}{4}=0\ \Rightarrow\ 5\varphi_A-2\varphi_B-\varphi_C=48\\
B:&\ \tfrac{\varphi_B-\varphi_A}{2}+\tfrac{\varphi_B-\varphi_C}{1}+\tfrac{\varphi_B}{2}=0\ \Rightarrow\ -\varphi_A+4\varphi_B-2\varphi_C=0\\
C:&\ \tfrac{\varphi_C-\varphi_A}{4}+\tfrac{\varphi_C-\varphi_B}{1}+\tfrac{\varphi_C-14}{2}=0\ \Rightarrow\ -\varphi_A-4\varphi_B+7\varphi_C=28
\end{aligned}
\]
\[
\boxed{\varphi_A=16\uu V,\quad\varphi_B=10\uu V,\quad\varphi_C=12\uu V.}
\]
\textbf{Corrientes}: $I_1=\frac{24-16}{2}=4\uu A$; $I_2=\frac{6}{2}=3\uu A$; $I_3=\frac44=1\uu A$; $I_4=\frac21=2\uu A$; $I_5=\frac{10}{2}=5\uu A$; $I_6=\frac{14-12}{2}=1\uu A$. Todas positivas: los sentidos supuestos eran los reales. Las tres ecuaciones de malla se cumplen: $8+4-2=10$, $-6+2+4=0$, $2+10+2=14$ \checkmark.

\textbf{Balance}: fuentes: $24(4)+14(1)=110\uu W$. Joule: $32+18+4+4+50+2=110\uu W$ \checkmark.

\sol{12.- Dos fuentes reales en paralelo (teorema de Thévenin)}
Sea $U$ la tensión común sobre $R$. Nodo superior: $\dfrac{\fem_1-U}{r_1}+\dfrac{\fem_2-U}{r_2}=\dfrac UR$, es decir,
\[
U\Big(\frac1{r_1}+\frac1{r_2}+\frac1R\Big)=\frac{\fem_1}{r_1}+\frac{\fem_2}{r_2}.
\]
Para una única fuente $(\fem_{eq},r_{eq})$: $U\Big(\dfrac1{r_{eq}}+\dfrac1R\Big)=\dfrac{\fem_{eq}}{r_{eq}}$. Como debe valer para todo $R$, se identifican los coeficientes:
\[
\frac1{r_{eq}}=\frac1{r_1}+\frac1{r_2}\Rightarrow\boxed{r_{eq}=\frac{r_1r_2}{r_1+r_2}},\qquad\fem_{eq}=r_{eq}\Big(\frac{\fem_1}{r_1}+\frac{\fem_2}{r_2}\Big)=\boxed{\frac{\fem_1r_2+\fem_2r_1}{r_1+r_2}}.
\]
\textbf{$I_2=0$}: $U=\fem_2$, y la fuente 1 alimenta sola a $R$: $\fem_2=\dfrac{\fem_1R}{R+r_1}\Rightarrow\boxed{R=\dfrac{\fem_2r_1}{\fem_1-\fem_2}}$ (solo es posible si $\fem_1>\fem_2$).

\sol{13.- Escalera R--2R}
\begin{center}\circEscalera{$R$}{$2R$}\end{center}
Autosemejanza: $R_\infty=R+\dfrac{2R\,R_\infty}{2R+R_\infty}\Rightarrow R_\infty^2-RR_\infty-2R^2=0\Rightarrow\boxed{R_\infty=2R}$.

\textbf{Tensiones}: desde el nodo 1 hacia la derecha se ve $2R\,\|\,R_\infty=2R\|2R=R$. El nodo 1 es un divisor entre la $R$ en serie y esa $R$: $U_1=U_0/2$. Lo mismo ocurre en cada celda, así que
\[
\boxed{U_n=\frac{U_0}{2^n}}.
\]
\textbf{Aplicación}: la red R--2R es la base de los convertidores digital--analógico (DAC). Cada bit conmuta una corriente de peso $2^{-n}$ usando solo dos valores de resistencia, fáciles de fabricar con precisión.

\sol{14.- Cubo: diagonal de cara y arista}
Vértices $(x,y,z)\in\{0,1\}^3$.

\textbf{(a) Diagonal de cara}, entre $A=(0,0,0)$ y $B=(1,1,0)$, con $\varphi_A=U$ y $\varphi_B=0$. La reflexión $x\leftrightarrow y$ fija $A$ y $B$ e intercambia $(1,0,0)\leftrightarrow(0,1,0)$. La simetría $(x,y)\to(1-x,1-y)$ intercambia $A\leftrightarrow B$ y lleva $(1,0,0)\to(0,1,0)$, así que esos dos vértices cumplen $\varphi=U-\varphi$: valen $U/2$. Igual ocurre con $(1,0,1)$ y $(0,1,1)$. Sea $p=\varphi(0,0,1)$; su nodo da $(p-U)+2(p-U/2)=0\Rightarrow p=2U/3$. Corriente que sale de $A$:
\[
I=\frac{U/2}{R}+\frac{U/2}{R}+\frac{U/3}{R}=\frac{4U}{3R}\ \Rightarrow\ \boxed{R_{eq}=\frac34R}.
\]
\textbf{(b) Arista}, entre $A=(0,0,0)$ y $B=(1,0,0)$, con $\varphi_A=+V$ y $\varphi_B=-V$. La antisimetría $x\to1-x$ da $\varphi(1,y,z)=-\varphi(0,y,z)$ y la simetría $y\leftrightarrow z$ da $\varphi(0,1,0)=\varphi(0,0,1)=p$; sea $\varphi(0,1,1)=s$.
Nodo $(0,1,0)$: $(p-V)+(p-(-p))+(p-s)=0\Rightarrow4p-s=V$. Nodo $(0,1,1)$: $2(s-p)+(s-(-s))=0\Rightarrow s=p/2$. Por lo tanto $p=2V/7$.
\[
I=\frac{2V}{R}+2\,\frac{V-2V/7}{R}=\frac{24V}{7R},\qquad R_{eq}=\frac{2V}{I}=\boxed{\frac{7}{12}R}.
\]
Orden lógico: $R_{\text{arista}}<R_{\text{cara}}<R_{\text{opuestos}}$, es decir, $7/12<3/4<5/6$ \checkmark.

\sol{15.- Rueda hexagonal}
\begin{center}\circRueda\end{center}
\textbf{(a) Vértices 1 y 4} ($\varphi_1=+V$, $\varphi_4=-V$). Por antisimetría, $\varphi_O=0$; por la simetría respecto del eje 1--4, $\varphi_2=\varphi_6=a$ y $\varphi_3=\varphi_5=-a$. Nodo 2 (vecinos 1, 3 y $O$): $(a-V)+(a+a)+(a-0)=0\Rightarrow a=V/4$.
\[
I=2\frac{V-V/4}{R}+\frac{V}{R}=\frac{5V}{2R},\qquad R_{14}=\frac{2V}{I}=\boxed{\frac45R}.
\]
\textbf{(b) Centro $O$ y vértice 1} ($\varphi_O=U$, $\varphi_1=0$). Por la simetría respecto del eje $O$--1: $\varphi_2=\varphi_6=x$, $\varphi_3=\varphi_5=y$, $\varphi_4=z$.
\[
\begin{aligned}
\text{nodo 2: }& x+(x-y)+(x-U)=0\Rightarrow3x-y=U\\
\text{nodo 3: }& (y-x)+(y-z)+(y-U)=0\Rightarrow3y-x-z=U\\
\text{nodo 4: }& 2(z-y)+(z-U)=0\Rightarrow3z-2y=U
\end{aligned}
\quad\Rightarrow\quad x=\frac{11}{18}U,\ y=\frac56U,\ z=\frac89U.
\]
Corriente que llega a 1: $\dfrac UR+2\dfrac xR=\dfrac{20U}{9R}$, así que $\boxed{R_{O1}=\dfrac{9}{20}R}$.

%--------------------------------------------------------------
\subsection*{Grupo IV}
\sol{16.- 24 pilas: agrupación óptima}
Condición teórica: $mr/k=R$ con $mk=24$, es decir, $m^2=\dfrac{NR}{r}=\dfrac{24(1.5)}{0.50}=72$, $m=8.49$: no es entero ni divide a 24. Se comparan los divisores de 24:
\begin{center}
\begin{tabular}{lccccc}
\toprule
$m\times k$ & $4\times6$ & $6\times4$ & $8\times3$ & $12\times2$ & $24\times1$\\
\midrule
$I=\dfrac{1.5m}{1.5+0.5m/k}$ (A) & 3.27 & 4.00 & \textbf{4.24} & 4.00 & 2.67\\
\bottomrule
\end{tabular}
\end{center}
\textbf{Óptimo}: 3 filas de 8 pilas en serie. $\fem_{eq}=12\uu V$, $r_{eq}=8(0.5)/3=1.33\ohm$:
\[
I=\frac{12}{1.5+1.333}=\boxed{4.24\uu A},\qquad P_R=I^2R=\boxed{26.9\uu W}.
\]
La condición $r_{eq}=R$ no se cumple exactamente porque $m$ y $k$ deben ser enteros con $mk=N$; el máximo real es el divisor más cercano al valor teórico.

\sol{17.- Carga de un capacitor a través de una resistencia}
\textbf{Sin ecuaciones diferenciales}: la fuente ideal mueve la carga final $Q=C\fem$ a través de su fem constante, así que realiza el trabajo $W_\fem=Q\fem=C\fem^2$. El capacitor almacena $W_C=\tfrac12C\fem^2$. No hay otras formas de energía (sin radiación ni inductancia), por lo tanto
\[
Q_J=W_\fem-W_C=\boxed{\tfrac12C\fem^2},\quad\text{independiente de }R.
\]
\textbf{Evolución}: malla, $\fem=IR+q/C$ con $I=dq/dt$, es decir, $\dfrac{dq}{C\fem-q}=\dfrac{dt}{RC}$. Con $q(0)=0$:
\[
q(t)=C\fem\big(1-e^{-t/RC}\big),\qquad I(t)=\frac{\fem}{R}e^{-t/RC}.
\]
\textbf{Verificación}: $\displaystyle\int_0^\infty I^2R\,dt=\frac{\fem^2}{R}\cdot\frac{RC}{2}=\frac12C\fem^2$ \checkmark.

\textbf{Numérico}: $Q=1.20\uu{mC}$; $W_\fem=14.4\uu{mJ}$; $W_C=7.20\uu{mJ}$; $Q_J=7.20\uu{mJ}$. La mitad de la energía entregada se pierde siempre, sea cual sea $R$.

\sol{18.- Dos interruptores cruzados}
\begin{center}\circInterruptoresDos\end{center}
\textbf{Ambos abiertos}: $3R\,\|\,3R=\dfrac{3R}{2}$.

\textbf{(a) Solo $K_1$} ($c\equiv f=X$): entre $a$ y $X$, la resistencia $a$--$c$ ($R$) está en paralelo con el camino $a$--$e$--$f$ ($2R$): $\dfrac{2R}{3}$. Entre $X$ y $b$, la resistencia $f$--$b$ ($R$) está en paralelo con $c$--$d$--$b$ ($2R$): $\dfrac{2R}{3}$. En total, $\dfrac{4R}{3}$:
\[
\Delta R_a=\frac{4R}{3}-\frac{3R}{2}=\boxed{-\frac R6}.
\]
\textbf{(b) $K_1$ y $K_2$} ($X=c\equiv f$, $Y=d\equiv e$): $a$--$X$ ($R$), $a$--$Y$ ($R$), $X$--$b$ ($R$), $Y$--$b$ ($R$), y entre $X$ e $Y$ las resistencias $c$--$d$ y $e$--$f$ en paralelo ($R/2$). La red es simétrica al intercambiar $X\leftrightarrow Y$, así que $\varphi_X=\varphi_Y$ y por $R/2$ no pasa corriente: $R_{ab}=(2R)\|(2R)=R$.
\[
\Delta R_b=R-\frac{3R}{2}=\boxed{-\frac R2}.
\]
\textbf{Observación}: cerrar $K_2$ después de $K_1$ reduce la resistencia en otros $R/3$.

\sol{19.- Capacitor en una red con fuentes}
\begin{center}\circSdosCap\end{center}
\textbf{Régimen estacionario}: por la rama central ($\fem_3$, $C$) no pasa corriente. Por la malla exterior circula $I$ en sentido horario (sube por $\fem_1$ y baja por $\fem_2$):
\[
I=\frac{\fem_1-\fem_2}{R_1+R_2}=\frac{12}{4}=\boxed{3.0\uu A},\qquad\varphi_A-\varphi_B=\fem_1-IR_1=20-9=11\uu V\ (=8+3\cdot1\ \checkmark).
\]
($\fem_2$ recibe energía: la corriente entra por su borne $+$, a razón de $24\uu W$.)

\textbf{Rama central}: de $B$ a $A$ se sube $\fem_3$ y luego la tensión del capacitor $U_C=\varphi_{\text{placa sup.}}-\varphi_{\text{placa inf.}}$: $11=5+U_C\Rightarrow U_C=6.0\uu V$. Queda positiva la \textbf{placa unida a $A$}, con $Q=CU_C=\boxed{24\uu{\mu C}}$.

\textbf{Se corta la rama de $\fem_2$}: no queda ningún camino cerrado de corriente continua, así que $I=0$, $\varphi_A-\varphi_B=\fem_1=20\uu V$, $U_C=15\uu V$ y $Q'=\boxed{60\uu{\mu C}}$. La carga adicional de la placa superior, $60-24=\boxed{36\uu{\mu C}}$, llega a través de $R_1$ (sube por la rama de $\fem_1$).

\sol{20.- Capacitor cilíndrico con dos capas con pérdidas}
Resistencia de fuga: dos capas de cascarones en serie, con $dR=\dfrac{dr}{\sigma\,2\pi rH}$:
\[
R=\frac{1}{2\pi H}\Big[\frac{\ln(R_0/R_1)}{\sigma_1}+\frac{\ln(R_2/R_0)}{\sigma_2}\Big],\qquad\boxed{I=\frac UR=\frac{2\pi HU}{\frac{1}{\sigma_1}\ln\frac{R_0}{R_1}+\frac{1}{\sigma_2}\ln\frac{R_2}{R_0}}}.
\]
Campo: $j(r)=\dfrac{I}{2\pi rH}$ en ambas capas (continuidad), así que
\[
\vE_1(r)=\frac{I}{2\pi\sigma_1Hr}\hat r\ (R_1<r<R_0),\qquad\vE_2(r)=\frac{I}{2\pi\sigma_2Hr}\hat r\ (R_0<r<R_2).
\]
\textbf{Carga libre en $r=R_0$}: $\sigma_f=D_2-D_1=\varepsilon_2E_2-\varepsilon_1E_1=\boxed{\dfrac{I}{2\pi R_0H}\Big(\dfrac{\varepsilon_2}{\sigma_2}-\dfrac{\varepsilon_1}{\sigma_1}\Big)}$, con carga total $q_{R_0}=I\Big(\dfrac{\varepsilon_2}{\sigma_2}-\dfrac{\varepsilon_1}{\sigma_1}\Big)$.

\textbf{Armadura interior}: $q_1=D_1(R_1)\,2\pi R_1H=\boxed{\dfrac{\varepsilon_1I}{\sigma_1}}$; la exterior tiene $q_2=-\dfrac{\varepsilon_2I}{\sigma_2}$. Comprobación: $q_1+q_{R_0}+q_2=0$ \checkmark. Si $\varepsilon_1/\sigma_1=\varepsilon_2/\sigma_2$, la interfaz queda sin carga (igual que en el problema 8).

\textbf{Potencia disipada}: $P=UI=\boxed{U^2/R}$. También se obtiene integrando $\sigma E^2$ en el volumen de las dos capas.


\end{document}
